% Direct-authority diplomatic French for printed p.166.
\SourcePage{166}

\LayerTarget{tohoku:remark:3.4.2:3:p2}
faisceau structural de $E$ puisse se réduire au faisceau \textit{constant}
$G$, i. e. pour que sur $E$ existe une connexion holomorphe intégrable ; et
cette condition est suffisante quand $X$ est de dimension complexe 1 (car
alors $u$ est épijective). L'invariant $d(E)$ a été d'abord introduit par
Weil [21] ; une définition plus géométrique de Atiyah [1] prouve que
l'annulation de $d(E)$ est en tous cas nécessaire et suffisante pour
l'existence d'une connexion holomorphe (non nécessairement intégrable) sur
$E$.

\LayerTarget{tohoku:section:3.5}
\indent\textbf{3.5. La suite exacte relative à un sous-espace fermé.} Soit $Y$ un
sous-espace localement fermé (i. e. intersection d'une partie ouverte et d'une
partie fermée) de l'espace $X$. Pour tout faisceau abélien $F$ sur $Y$, il
existe un faisceau abélien et un seul sur $X$, dont la restriction à $Y$ soit
$F$ et dont la restriction à $\complement Y$ soit 0. Pour le voir, on est ramené aussitôt
au cas où $Y$ est ouvert ou fermé, où la vérification facile est faite dans
[4, XVII, prop. 1].
\LayerTarget{tohoku:def:extension-by-zero-locally-closed}
Ce faisceau sur $X$ sera noté $F^X$. $F\to F^X$ est un foncteur exact
$\mathbf C^Y\to \mathbf C^X$, de plus on a, si $Z\subset Y\subset X$ ($Z$
localement fermé dans $Y$ donc dans $X$), et si $F$ est un faisceau abélien
sur $Z$ :
\LayerTarget{tohoku:eq:extension-by-zero-transitivity}
$(F^Y)^X=F^X$.
\LayerTarget{tohoku:def:restriction-extension-FY}
Si maintenant $F$ est un faisceau abélien sur $X$, on pose
$F_X=(F\vert Y)^X$, c'est donc le faisceau sur $X$ caractérisé par la
condition que sa restriction à $Y$ est la même que celle de $F$, et sa
restriction à $\complement Y$ est nulle. $F_Y$ est un foncteur exact
$\mathbf C^X\to \mathbf C^X$, de plus on a encore une propriété de transitivité
\LayerTarget{tohoku:eq:restriction-extension-transitivity}
$(F_Y)_Z=F_Z$
si $Z\subset Y\subset X$ sont comme ci-dessus. Si on suppose $Y$ fermé, donc
$U=\complement Y$ ouvert, on a une suite exacte bien connue

\LayerTarget{tohoku:eq:closed-subspace-exact-sequence}
\[
0\to F_U\to F\to F_Y\to0
\]
pour tout faisceau abélien $F$ sur $X$. Rappelons que l'on a
$\Gamma(F_Y)=\Gamma(Y,F)=\Gamma(F\vert Y)$, tandis que $\Gamma(F_U)$
s'identifie au sous-groupe de $\Gamma(F)$ formé des sections dont le support
est contenu dans $U$.
\LayerTarget{tohoku:def:induced-support-antifilter}
Soit $\Phi$ un antifiltre de parties fermées de $X$ ; pour toute partie $Z$
de $X$, soit $\Phi_Z$ l'antifiltre ``induit'' formé des $A\in\Phi$ qui
sont contenus dans $Z$ (ne pas confondre avec la trace $\Phi\cap Z$ de $\Phi$
sur $Z$ !). On vérifie facilement que si $Z$ est localement fermé, et $\Phi$
paracompactifiante (cf. 3.3.) alors $\Phi_Z$ est aussi paracompactifiante. Si
$Z$ est localement fermé, $\Phi$ quelconque, on déduit facilement des formules
ci-dessus la formule plus générale suivante (valable en particulier si $Z$
est ouvert, ou fermé) :

\LayerTarget{tohoku:eq:sections-support-restriction}
\[
\Gamma_\Phi(X,F_Z)=\Gamma_{\Phi_Z}(Z,F\vert Z)
\]
valable pour tout faisceau abélien $F$ sur $X$ ; formule d'ailleurs équivalente
à la suivante :

\LayerTarget{tohoku:eq:sections-extension-by-zero}
\[
\Gamma_{\Phi_Z}(Z,G)=\Gamma_\Phi(X,G^X)
\]
valable pour tout faisceau abélien $G$ sur $Z$.
\LayerTarget{tohoku:construction:cohomology-support-comparison:p1}
Comme $G^X$ est un foncteur \textit{exact} en $G$, les $H^p_\Phi(X,G^X)$
forment un foncteur cohomologique sur $\mathbf C^Z$, et comme le foncteur cohomologique
universel $(H^p_{\Phi_Z}(Z,G))$ coïncide avec le premier en dimension 0, on en
conclut des \textit{homomorphismes canoniques}

\LayerTarget{tohoku:eq:canonical-support-comparison}
\[
H^p_{\Phi_Z}(Z,G)\to H^p_Z(X,G^X)
\]
(caractérisés par le fait de définir un homomorphisme de foncteurs coho-

\TohokuContinuation{167}
